\section{A clear path to complete Lean formalization}\label{sec:lean}

\subsection{The object-theory target}

The inspected ProveIt snapshot is \texttt{\pin}, using Lean~4.32.0.
Its \href{https://github.com/VladimirReshetnikov/ProveIt/tree/d851229c33ef395db8dc04baf28c6a291073f61e/Logic/PrimitiveRecursiveArithmetic/ConciseIPRA/Lean/ConciseIPRA}{\texttt{ConciseIPRA} sources}
already supply scoped syntax (\texttt{Syntax}), intuitionistic derivability
(\texttt{Deduction}), primitive-recursive function symbols and
quantifier-free induction (\texttt{PR}), and transport to the concise
native calculus (\texttt{NativeBiInterpretation}).
The required endpoint is a derivation \emph{inside that encoded theory}:
\begin{lstlisting}
open LeanProofs.ConciseIPRA
Prf PRAAx ([] : List (Fm PTm 0))
  (PuiseuxIPRA.PuiseuxSentence concretePrograms)
\end{lstlisting}
Here \texttt{concretePrograms} must contain the actual \texttt{PR} syntax
trees for input validation, packet construction, coefficient generation,
and the finite predicates of Section~\ref{sec:internal}.  Their decoding
theorems must establish the intended meanings; naming a Boolean program
\texttt{factorJet} does not establish its correctness.

The supplied \texttt{lean/PuiseuxTarget.lean} writes the closed formula
explicitly using \texttt{Fm.eq}, \texttt{Fm.and}, \texttt{Fm.imp}, and
\texttt{Fm.all}.  Its matrix asserts packet validity, the complete
denominator-cleared factorization identity, prefix compatibility, and
finite coefficient bounds; a second universally quantified matrix asserts
the rational tail estimate.  It also gives the type of the required
derivation and conditional transport of an already supplied derivation
through \texttt{pra\_iff\_native}.  This last transport does not construct
the missing derivation.

\subsection{Dependency contracts and implementation order}

The following sequence specifies concrete completion targets.  Every
algorithm operates on natural-number codes, with default outputs on
malformed data; valid-input proofs exclude those defaults.

\begin{enumerate}
\setlength{\itemsep}{3pt}
\item \textbf{Coded arithmetic.}
Implement signed integers, rational numbers with positive denominators,
finite lists, polynomial arrays, bounded scans and Boolean folds.
Prove decoding equations for addition, multiplication, comparison,
coefficient access and finite convolution.

\item \textbf{Exact algebraic coefficients.}
Formalize Section~\ref{sec:backend}: real-root isolation, algebraic
arithmetic, equality and order, resultant elimination and finite splitting.
The reused constructive FTA root-name construction requires its own
deep arithmetic derivation; the existing draft is not a compiled
FTA theorem.

\item \textbf{Squarefree preprocessing.}
Prove the degree-decreasing Euclidean algorithm and its product and
B\'ezout invariants.  Verify the multiplicities, common denominator \(C\),
monic polynomial factors \(F_m\), and the denominator-cleared identity
relating them to the original polynomial.

\item \textbf{Truncated polynomial algebra.}
Prove congruence preservation for multiplication, substitution,
ramification, translation and discriminants.  Prove the finite
Hasse--Taylor identity and guarded division by a monomial.
These identities determine the exact precision losses.

\item \textbf{Coprime Hensel recursion.}
Implement the correction
\(U_{j,k}=\operatorname{rem}(E_jR_k,C_j)\).
Prove the product identity through the next coefficient, uniqueness
at the current precision, and compatibility when the requested
precision is increased.

\item \textbf{Bounded Newton tree.}
Implement Section~\ref{sec:newton} as a bounded stack or \(d\) rounds,
with at most \(2d-1\) nodes.  Maintain degree accounting,
\(e\mid N\), the discriminant bound \(e\delta\), and the precision reserve.
The scan bound \(B\) and lookahead \(J(M)\) are part of the program.

\item \textbf{Finite discriminant separation.}
From the factorization modulo \(T^M\), prove the discriminant congruence,
find every first difference by a scan through \(\Delta\), and establish
\(2\sum_{i<j}w_{ij}=\Delta\).
Only equality of finite algebraic coefficients is tested.

\item \textbf{Simple polynomial leaves.}
Formalize Section~\ref{sec:bootstrap} with \(a_i\) through degree \(h\)
and \(M=2h+1\).  Prove \(s_i\le\Delta\), the divisibility defining \(H_i\),
\(H_i(0,0)=0\), and \(H_{i,U}(0,0)\ne0\).
All four conclusions are finite polynomial assertions.

\item \textbf{Coherent coefficients and complete factorization.}
Implement the simple coefficient recurrence and prove its finite
correctness and uniqueness.  Synthetic division and cancellation by
the recorded leading differences recover every factor \(F_m\).
Use residual congruences modulo \(T^{L+\Delta+1}\) when establishing
factorization through degree \(L\).

\item \textbf{Rational majorants and evaluation.}
Formalize Section~\ref{sec:analytic} with truncated positive Picard
polynomials, rational coefficient majorants and finite weighted sums.
Add unit inversion and convolution to obtain the common radius and
Laurent tail bound, including the supplied positive lower bound
for a nonzero argument.

\item \textbf{Internal proof assembly.}
Construct the programs from \texttt{PR.comp} and \texttt{PR.prec}.
For each loop, use \texttt{PRA.qfInd} on its finite Boolean invariant
with the horizon held as a parameter.  Universally generalize afterward
and assemble the displayed \texttt{Prf PRAAx} endpoint.

\item \textbf{Semantic interpretation and audit.}
Interpret the coherent coefficient and tail statements as convergent
Puiseux expansions.  Transfer the object derivation to the native
calculus if desired, inspect its type and \texttt{\#print axioms},
and rebuild its complete dependency closure from the pinned sources.
\end{enumerate}

\subsection{Current status and the remaining work}

No Lean executable was available in this run.  The supplied Lean source
is therefore an \emph{uncompiled interface and outline}, including its
small structural proof terms.  It contains no omitted proof term or
postulated Puiseux theorem, and supplies no inhabitant of the final
derivation type.  The exact-algebraic backend, quantitative FTA dependency,
Newton--Hensel invariants, majorant bounds and their object-theory
derivations remain substantial formalization tasks.

The accompanying \texttt{lean/LEAN\_PLAN.md} expands the contracts above,
including the exact original-polynomial identity checked by
\texttt{factorJet}.  The reference computations provide finite evidence
for executable components.  A conventional Lean theorem about
\texttt{Complex} or \texttt{PowerSeries} may help organize the mathematics,
but it does not replace the encoded \(\ipra\) derivation required here.

